% G2-5 Alur satu iterasi AdaBoost.R2.
% Sumber: diadaptasi dari drucker1997 (hlm. PDF 2--3); batas jumlah pohon M
% (`n_estimators`) dari sklearnadaboost.
% Satu kolom, dua jalur berhenti digabung di sisi kanan, jalur ulang di kiri.
\begin{tikzpicture}[x=1mm,y=-1mm,font=\footnotesize,
    langkah/.style={draw,rectangle,text width=52mm,align=center,
      inner sep=2.5pt,minimum height=11mm},
    kep/.style={draw,diamond,aspect=2.8,inner sep=0pt,text width=22mm,
      align=center},
    hasil/.style={draw,rounded rectangle,text width=52mm,align=center,
      inner sep=2.5pt,minimum height=12mm},
    panah/.style={-{Stealth[length=2mm]}}]

  % --- rantai utama, satu kolom, arah ke bawah ---
  \node[langkah] (s1) at (40,0)
    {(1) bobot sampel $w_i$ menjadi probabilitas
     $p_i = w_i / \sum_j w_j$};
  \node[langkah] (s2) at (40,19)
    {(2) ambil sampel ulang dengan pengembalian menurut $p_i$};
  \node[langkah] (s3) at (40,38)
    {(3) latih pohon regresi ke-$m$};
  \node[langkah] (s4) at (40,57)
    {(4) hitung \textit{loss} $L_i$ dan $\bar{L} = \sum_i L_i\,p_i$};
  \node[kep]     (d1) at (40,78) {$\bar{L} \ge 0{,}5$?};
  \node[langkah] (s5) at (40,101)
    {(5) hitung $\beta_m = \bar{L}/(1-\bar{L})$ lalu perbarui
     $w_i \leftarrow w_i\,\beta_m^{\,1-L_i}$};
  \node[kep]     (d2) at (40,126) {$m = M$?};
  \node[hasil]   (fin) at (40,152)
    {\textbf{ramalan akhir}\\median berbobot $f_m(x)$ dengan bobot
     $\log(1/\beta_m)$};

  \draw[panah] (s1) -- (s2);
  \draw[panah] (s2) -- (s3);
  \draw[panah] (s3) -- (s4);
  \draw[panah] (s4) -- (d1);
  \draw[panah] (d1) -- (s5);
  \node[anchor=west] at (41,89) {tidak};
  \draw[panah] (s5) -- (d2);

  % --- dua jalur "ya" digabung di sisi kanan, lalu masuk ramalan akhir ---
  \draw (d1.east) -- (100,78);
  \node[anchor=south west] at (62,77) {ya};
  \draw (d2.east) -- (100,126);
  \node[anchor=south west] at (62,125) {ya};
  \fill (100,126) circle (0.8mm);
  \draw (100,78) -- (100,152);
  \draw[panah] (100,152) -- (fin.east);

  % --- jalur ulang di sisi kiri ---
  \draw (d2.west) -- (-14,126);
  \draw (-14,126) -- (-14,0);
  \draw[panah] (-14,0) -- (s1.west);
  \node[anchor=north west,align=left] at (-13,128) {tidak, $m \leftarrow m+1$};
\end{tikzpicture}
